\chapter{Preferences and Utility}\label{ch:b1:preferences-and-utility}

Between 1996 and 2024 the prices that households in the euro area paid
moved far apart. Electricity cost two and a half times as much at the end
as at the start; televisions, cameras and other audio-visual equipment cost
less than a fifth as much. A household that spent the same income in both
years could not have bought the same basket, and would not have wanted to.
To say what it would do instead, a model needs a description of what the
household wants that does not change when prices change. This chapter
builds that description: a preference ordering over baskets, the axioms
that make it usable, and the \omterm{def:b1:preferences-and-utility:utility}{utility functions} that represent it. The
next chapter puts it to work against a budget.

\begin{openingchart}{Consumer prices in the euro area, five categories and all
items, 1996--2024, annual average indices with 2015 = 100, on a logarithmic
axis.}{Eurostat, harmonised index of consumer prices, annual data
(\code{prc_hicp_aind}), euro area}
\begin{tikzpicture}
\begin{semilogyaxis}[omtimesten, xlabel={Year}, ylabel={Index, 2015 = 100},
  xmin=1996, xmax=2024, ymin=50, ymax=500, ytick={50,100,200,500},
  yticklabels={50,100,200,500}, width=0.72\linewidth,
  height=0.5\linewidth, legend pos=outer north east]
  \addplot+[very thick, black] table[x=year, y=all] {figdata/bachelor-1/04-preferences-and-utility/out/hicp_categories.csv};
  \addplot table[x=year, y=electricity] {figdata/bachelor-1/04-preferences-and-utility/out/hicp_categories.csv};
  \addplot table[x=year, y=bread] {figdata/bachelor-1/04-preferences-and-utility/out/hicp_categories.csv};
  \addplot table[x=year, y=transport] {figdata/bachelor-1/04-preferences-and-utility/out/hicp_categories.csv};
  \addplot table[x=year, y=rent] {figdata/bachelor-1/04-preferences-and-utility/out/hicp_categories.csv};
  \addplot table[x=year, y=audiovisual] {figdata/bachelor-1/04-preferences-and-utility/out/hicp_categories.csv};
  \legend{All items, Electricity, Bread and cereals, Transport services, Rents, Audio-visual}
\end{semilogyaxis}
\end{tikzpicture}
\end{openingchart}

All six series are set to 100 in 2015, so the chart compares changes, not
levels. From 1996 to 2024 the index of all items rose by a factor of 1.757.
Electricity rose by a factor of 2.547, bread and cereals by 1.932, rents by
1.603, and audio-visual equipment fell to 0.188 of its 1996
level. % figdata: hicp_ratios
Relative to audio-visual equipment, electricity became
$2.547/0.188 = 13.5$ times dearer. % figdata: hicp_ratios
A household that kept its 1996 basket through 2024 would have bought too
much of what had become dear and too little of what had become cheap. How
much it moves away from that basket depends on how willing it is to
substitute one good for another, and that is a property of its
preferences, which the chapter now defines.

\begin{recall}
Two goods are substitutes when a rise in the price of one raises the demand
for the other, and complements when it lowers it. The law of demand says
that the quantity demanded falls when the good's own price rises. A
household's budget is the income it can spend.
\end{recall}

\section{Preference relations}\label{sec:b1:preferences-and-utility:relations}

A \emph{bundle} is a list of quantities $x = (x_1, x_2)$ of two goods, with
$x_1, x_2 \ge 0$; everything in this chapter extends to $n$ goods. A
consumer compares bundles two at a time.

\begin{definition}[Preference relation]\label{def:b1:preferences-and-utility:preference}
\index{preference relation}\index{weak preference}\index{strict preference}\index{indifference}
A \emph{preference relation} $\succsim$ on bundles reads $x \succsim y$: ``$x$
is at least as good as $y$'' for the consumer (\emph{weak preference}). From
it, $x \succ y$ (\emph{strict preference}) means $x \succsim y$ and not
$y \succsim x$, and $x \sim y$ (\emph{indifference}) means $x \succsim y$ and
$y \succsim x$.
\end{definition}

\begin{definition}[Rational preferences]\label{def:b1:preferences-and-utility:rational}
\index{completeness (of preferences)}\index{transitivity (of preferences)}\index{rational preferences}
A \omterm{def:b1:preferences-and-utility:preference}{preference relation} is \emph{complete} if for any two bundles $x$, $y$,
$x \succsim y$ or $y \succsim x$ (or both). It is \emph{transitive} if
$x \succsim y$ and $y \succsim z$ imply $x \succsim z$. A complete and transitive
\omterm{def:b1:preferences-and-utility:preference}{preference relation} is called \emph{rational}.
\end{definition}

Completeness rules out a consumer who cannot compare two bundles at all.
Transitivity rules out cycles: someone who prefers tea to coffee, coffee to
juice and juice to tea would pay a little to swap tea for juice, a little
more to swap juice for coffee, a little more to swap coffee for tea, and end
where they started, poorer. The word ``\omterm{def:b1:preferences-and-utility:rational}{rational}'' means only these two
properties, nothing about whether the tastes are wise.

\begin{definition}[Monotone preferences]\label{def:b1:preferences-and-utility:monotone}
\index{monotonicity (of preferences)}
Preferences are \emph{monotone} if more of every good is better: whenever
$x_1 \ge y_1$ and $x_2 \ge y_2$ with at least one inequality strict,
$x \succ y$.
\end{definition}

Monotonicity rules out satiation in the goods modelled: no bundle is so
large that more of it is unwelcome. It is an assumption about the goods
chosen for the model (a good that has become a nuisance is modelled as its
removal, which is a good).

\begin{definition}[Convex preferences]\label{def:b1:preferences-and-utility:convex}
\index{convex preferences}\index{better-than set}
For a bundle $y$, the \emph{better-than set} is $\{x : x \succsim y\}$.
Preferences are \emph{convex} if every better-than set is convex: if
$x \succsim y$ and $z \succsim y$, then $tx + (1 - t)z \succsim y$ for every
$t \in [0, 1]$. They are \emph{strictly convex} if, moreover,
$tx + (1 - t)z \succ y$ whenever $x \neq z$ and $0 < t < 1$.
\end{definition}

Convexity says that averages are at least as good as extremes: a consumer
indifferent between four litres of milk with two loaves and two litres
with four loaves likes three of each at least as much. It rules out a taste
for specialising in one good.

\section{Indifference curves and the marginal rate of substitution}\label{sec:b1:preferences-and-utility:curves}

\begin{definition}[Indifference curve]\label{def:b1:preferences-and-utility:indifference-curve}
\index{indifference curve}
The \emph{\omterm{def:b1:preferences-and-utility:preference}{indifference} curve} through a bundle $y$ is the set of bundles
indifferent to it, $\{x : x \sim y\}$.
\end{definition}

\begin{omfigure}
\begin{tikzpicture}
\begin{axis}[omrecord, grid=none, xlabel={Good 1, $x_1$}, ylabel={Good 2, $x_2$},
  xmin=0, xmax=8, ymin=0, ymax=8, width=0.5\linewidth, height=0.42\linewidth, samples=80]
  \fill[omDef!14] (axis cs:0,0) rectangle (axis cs:8,8);
  \addplot[fill=white, draw=none, domain=0.01:8] {min(9/x, 8)} \closedcycle;
  \addplot[guide] coordinates {(1.5,6) (6,1.5)};
  \addplot[demand, domain=1.125:8] {9/x};
  \addplot[only marks, mark=*, black] coordinates {(3,3) (1.5,6) (6,1.5)};
  \node[font=\footnotesize, anchor=south west] at (axis cs:3,3) {$y$};
  \node[font=\footnotesize] at (axis cs:5.6,5.6) {better than $y$};
\end{axis}
\end{tikzpicture}

{\small The \omterm{def:b1:preferences-and-utility:indifference-curve}{indifference curve} through $y = (3, 3)$ (blue) and the
\omterm{def:b1:preferences-and-utility:convex}{better-than set} above it (shaded). Convexity: the segment joining
$(1.5, 6)$ and $(6, 1.5)$ lies in the \omterm{def:b1:preferences-and-utility:convex}{better-than set}.}
\end{omfigure}

\begin{proposition}[Indifference curves do not cross]\label{prop:b1:preferences-and-utility:no-cross}
\index{indifference curve!no crossing}
If preferences are \omterm{def:b1:preferences-and-utility:rational}{rational}, two \omterm{def:b1:preferences-and-utility:indifference-curve}{indifference curves} are either the same
curve or have no point in common. If preferences are also \omterm{def:b1:preferences-and-utility:monotone}{monotone}, no two
bundles on one \omterm{def:b1:preferences-and-utility:indifference-curve}{indifference curve} have one bundle with more of both goods,
so every \omterm{def:b1:preferences-and-utility:indifference-curve}{indifference curve} slopes downward.
\end{proposition}

\begin{proof}
Let the curves through $y$ and $z$ share a bundle $x$: $x \sim y$ and
$x \sim z$. By transitivity applied in both directions, $y \sim z$, and then
every bundle indifferent to $y$ is indifferent to $z$ and conversely (again
by transitivity): the two curves are the same set. For the second part, if
$x$ and $x'$ lay on one curve with $x'$ larger in both goods, monotonicity
would give $x' \succ x$, contradicting $x' \sim x$.
\end{proof}

The slope of an \omterm{def:b1:preferences-and-utility:indifference-curve}{indifference curve} is the exchange rate between the two
goods at which the consumer is willing to trade.

\begin{definition}[Marginal rate of substitution]\label{def:b1:preferences-and-utility:mrs}
\index{marginal rate of substitution}\index{MRS}
At a bundle where the \omterm{def:b1:preferences-and-utility:indifference-curve}{indifference curve} is a differentiable graph
$x_2(x_1)$, the \emph{marginal rate of substitution} of good 1 for good 2 is
$\mathrm{MRS} = -\dd x_2/\dd x_1$ along the curve: the amount of good 2 the
consumer would give up for one more unit of good 1, in the limit of small
trades. Its unit is units of good 2 per unit of good 1.
\end{definition}

With strictly \omterm{def:b1:preferences-and-utility:convex}{convex} and \omterm{def:b1:preferences-and-utility:monotone}{monotone} preferences the \omterm{def:b1:preferences-and-utility:mrs}{MRS} falls as the consumer
moves down an \omterm{def:b1:preferences-and-utility:indifference-curve}{indifference curve}: the more of good 1 the consumer holds,
the less good 2 they will give up for another unit of it. This is
\emph{diminishing \omterm{def:b1:preferences-and-utility:mrs}{MRS}}; geometrically, the curve is \omterm{def:b1:preferences-and-utility:convex}{convex} towards the
origin.

\section{Utility functions}\label{sec:b1:preferences-and-utility:utility}

Comparing bundles two at a time is cumbersome. A function that gives every
bundle a number, larger for better bundles, turns comparisons into
arithmetic and optimisation into calculus.

\begin{definition}[Utility function]\label{def:b1:preferences-and-utility:utility}
\index{utility function}\index{ordinal utility}
A \emph{utility function} represents the \omterm{def:b1:preferences-and-utility:preference}{preference relation} $\succsim$ if,
for all bundles, $u(x) \ge u(y)$ if and only if $x \succsim y$. A utility
function in this sense is \emph{ordinal}: only the order of its values
carries meaning, not their differences or ratios.
\end{definition}

A \omterm{def:b1:preferences-and-utility:preference}{preference relation} that has a utility representation is \omterm{def:b1:preferences-and-utility:rational}{rational}, since
real numbers are completely and transitively ordered. The converse needs a
little more: the preferences must be continuous (no sudden reversal for a
small change in a bundle). One \omterm{def:b1:preferences-and-utility:rational}{rational} ordering fails it: the
\emph{lexicographic} order that compares good 1 first and looks at good 2
only to break ties. Each of its \omterm{def:b1:preferences-and-utility:indifference-curve}{indifference curves} is a single point, and
there are too many such points for the real numbers to rank them all; it
has no utility representation (a result of set theory, admitted here).

\begin{proposition}[Utility is ordinal]\label{prop:b1:preferences-and-utility:ordinal}
\index{utility function!increasing transformation}
If $u$ represents $\succsim$ and $f$ is a strictly increasing function on the
values of $u$, then $v = f \circ u$ represents $\succsim$. If $u$ and $f$ are
differentiable with $f' > 0$, then $v$ and $u$ give the same \omterm{def:b1:preferences-and-utility:mrs}{marginal rate
of substitution}:
\[
\mathrm{MRS} = \frac{\partial u/\partial x_1}{\partial u/\partial x_2} = \frac{\partial v/\partial x_1}{\partial v/\partial x_2}.
\]
\end{proposition}

\begin{proof}
$v(x) \ge v(y)$ if and only if $f(u(x)) \ge f(u(y))$, if and only if
$u(x) \ge u(y)$, since $f$ is strictly increasing. Along an \omterm{def:b1:preferences-and-utility:indifference-curve}{indifference
curve} $u(x_1, x_2) = c$, \cref{prop:b1:constrained-optimisation:level-slope}
gives $\dd x_2/\dd x_1 = -u_1/u_2$, so $\mathrm{MRS} = u_1/u_2$. For $v$,
$v_i = f'(u)u_i$, and the factor $f'(u) > 0$ cancels in $v_1/v_2$.
\end{proof}

The proposition is the reason the book never compares ``amounts of
utility'' between two people or between two bundles of one person. What a
\omterm{def:b1:preferences-and-utility:utility}{utility function} says about choices --- which bundle is chosen, how
choices respond to prices --- depends only on the \omterm{def:b1:preferences-and-utility:mrs}{MRS}, and the \omterm{def:b1:preferences-and-utility:mrs}{MRS} is the
same for every representation.

\section{The standard forms}\label{sec:b1:preferences-and-utility:forms}

Four families of \omterm{def:b1:preferences-and-utility:utility}{utility functions} do most of the work in the book.

\begin{definition}[Cobb--Douglas utility]\label{def:b1:preferences-and-utility:cobb-douglas}
\index{Cobb--Douglas utility}
A \emph{Cobb--Douglas utility} is $u(x_1, x_2) = x_1^{a}x_2^{1-a}$ with
$0 < a < 1$. Its \omterm{def:b1:preferences-and-utility:mrs}{MRS} is $\dfrac{a}{1 - a}\,\dfrac{x_2}{x_1}$.
\end{definition}

\begin{definition}[Perfect substitutes and perfect complements]\label{def:b1:preferences-and-utility:perfect}
\index{perfect substitutes}\index{perfect complements}
Goods are \emph{perfect substitutes} if preferences are represented by
$u = a x_1 + b x_2$ with $a, b > 0$: the \omterm{def:b1:preferences-and-utility:mrs}{MRS} is the constant $a/b$, and the
\omterm{def:b1:preferences-and-utility:indifference-curve}{indifference curves} are straight lines. Goods are \emph{perfect
complements} if preferences are represented by
$u = \min(x_1/a, x_2/b)$: they are consumed in the fixed proportion
$x_1 : x_2 = a : b$, and the \omterm{def:b1:preferences-and-utility:indifference-curve}{indifference curves} are right angles with their
corners on the ray of that proportion.
\end{definition}

\begin{definition}[Quasi-linear utility]\label{def:b1:preferences-and-utility:quasi-linear}
\index{quasi-linear utility}
A \emph{quasi-linear utility} is $u(x_1, x_2) = v(x_1) + x_2$ with $v$
increasing and concave. Its \omterm{def:b1:preferences-and-utility:mrs}{MRS} is $v'(x_1)$, which does not depend on $x_2$.
\end{definition}

\begin{omfigure}
\begin{tikzpicture}
\begin{groupplot}[group style={group size=2 by 2, horizontal sep=14mm, vertical sep=24mm},
  omrecord, grid=none, width=0.42\linewidth, height=0.32\linewidth,
  xmin=0, xmax=10, ymin=0, ymax=10, samples=60, xlabel={$x_1$}, ylabel={$x_2$}]
\nextgroupplot[title={\small Cobb--Douglas $\sqrt{x_1x_2}$}, width=0.42\linewidth]
  \addplot[demand, domain=0.4:10] {4/x};
  \addplot[demand, domain=0.9:10] {9/x};
  \addplot[demand, domain=1.6:10] {16/x};
\nextgroupplot[title={\small Perfect substitutes $x_1 + 2x_2$}, width=0.42\linewidth]
  \addplot[demand, domain=0:4] {(4 - x)/2};
  \addplot[demand, domain=0:8] {(8 - x)/2};
  \addplot[demand, domain=0:10] {(12 - x)/2};
\nextgroupplot[title={\small Perfect complements $\min(x_1, x_2)$}, width=0.42\linewidth]
  \addplot[demand] coordinates {(2,10) (2,2) (10,2)};
  \addplot[demand] coordinates {(4,10) (4,4) (10,4)};
  \addplot[demand] coordinates {(6,10) (6,6) (10,6)};
  \addplot[guide, domain=0:10] {x};
\nextgroupplot[title={\small Quasi-linear $4\sqrt{x_1} + x_2$}, width=0.42\linewidth]
  \addplot[demand, domain=0:10] {8 - 4*sqrt(x)};
  \addplot[demand, domain=0:10] {11 - 4*sqrt(x)};
  \addplot[demand, domain=0:10] {14 - 4*sqrt(x)};
\end{groupplot}
\end{tikzpicture}

{\small Three \omterm{def:b1:preferences-and-utility:indifference-curve}{indifference curves} of each standard form. Cobb--Douglas:
curves \omterm{def:b1:preferences-and-utility:convex}{convex} towards the origin. Substitutes: parallel lines. Complements:
right angles on the ray $x_2 = x_1$ (dotted). Quasi-linear: vertical shifts
of one curve, so the \omterm{def:b1:preferences-and-utility:mrs}{MRS} at a given $x_1$ is the same on every curve.}
\end{omfigure}

Cobb--Douglas preferences are the workhorse: smooth, strictly \omterm{def:b1:preferences-and-utility:convex}{convex},
with an \omterm{def:b1:preferences-and-utility:mrs}{MRS} that depends only on the ratio $x_2/x_1$, so that scaling a
bundle up leaves the \omterm{def:b1:preferences-and-utility:mrs}{MRS} unchanged. \omterm{def:b1:preferences-and-utility:perfect}{Perfect substitutes} model two brands
of the same thing. \omterm{def:b1:preferences-and-utility:perfect}{Perfect complements} model goods used together in fixed
amounts: left and right shoes, or flour and yeast in one recipe.
Quasi-linear preferences model a small good next to ``all other spending''
$x_2$, measured in money: the \omterm{def:b1:preferences-and-utility:mrs}{MRS} $v'(x_1)$ is then a willingness to pay in
money for one more unit of good 1, whatever the consumer's wealth, which is
what makes them the basis of consumer surplus (\cref{ch:b1:consumer-welfare}).

\begin{measure}[Relative prices from indices]
A price index with 2015 = 100 says how a price changed, not what it is.
The opening chart shows that electricity rose by a factor of 2.547 and
audio-visual equipment fell by a factor of 0.188 between 1996 and 2024; % figdata: hicp_ratios
it does not say that a kilowatt-hour costs more or less than a television.
The ratio of two indices is the change in a relative price: here
$2.547/0.188 = 13.5$. % figdata: hicp_ratios
The change in a relative price is what a model of choice needs, because an
\omterm{def:b1:preferences-and-utility:mrs}{MRS} is compared with a price ratio, and a price ratio that changes by a
factor of 13.5 changes by that factor whatever the units of the two goods.
An index of a category is itself an average over many products (and, for
equipment whose quality improves, an attempt to price the same quality
across years), so its ratio to another index is the relative price of two
averages, not of two goods.
\end{measure}

\section{Testing the axioms}\label{sec:b1:preferences-and-utility:testing}

Preferences are not observed; choices are. The axioms of this chapter have
testable consequences only together with a model of choice: if a consumer
with \omterm{def:b1:preferences-and-utility:rational}{rational}, \omterm{def:b1:preferences-and-utility:monotone}{monotone} preferences chooses the best bundle she can afford,
her choices across different budgets must be consistent with one another.
\Cref{ch:b1:consumer-choice} states the condition exactly (the generalised
axiom of revealed preference). Its idea is simple: if she chose $x$ when $y$
was affordable, she has revealed $x \succsim y$; a chain of such revelations
must never lead back to a bundle strictly worse than where it started.

\begin{model}{Rational preferences}
\assumes A consumer's preferences over bundles are complete, transitive,
\omterm{def:b1:preferences-and-utility:monotone}{monotone} and \omterm{def:b1:preferences-and-utility:convex}{convex}, and stable over the period studied.
\predicts Choices from different budgets reveal no cycle; \omterm{def:b1:preferences-and-utility:indifference-curve}{indifference
curves} slope down, do not cross and have a diminishing \omterm{def:b1:preferences-and-utility:mrs}{MRS}; a consumer
facing a change in relative prices moves along an \omterm{def:b1:preferences-and-utility:indifference-curve}{indifference curve}
towards the good that became cheaper.
\tested In a laboratory experiment with 93 subjects, each choosing from 50
budgets, 80 of the 93 scored above 0.90 on Afriat's efficiency index,
against 12 of 25,000 simulated random choosers (see the record box). % ledger: b1:choi2007-ccei
\breaks When the way the options are presented or a default option changes
choices that should be the same (the Year~2 volume's chapter on behavioural
evidence), and for goods whose consumption changes the preferences
themselves, such as addictive goods.
\end{model}

\begin{record}[Consistent choices in the laboratory]
Choi, Fisman, Gale and Kariv (2007; the counts below are from the
authors' 2006 working-paper version) gave each of 93 subjects 50 budget
lines, drawn at random, and asked them to divide tokens between two
accounts; one account would pay off at random. % ledger: b1:choi2007-ccei
A subject whose choices satisfy the axioms of revealed preference exactly
has an efficiency index of 1; an index of 0.9 means that the budgets would
have to be shrunk by $10\%$ to remove every inconsistency. 80 of the 93
subjects scored above 0.90, and 75 above 0.95. % ledger: b1:choi2007-ccei
Of 25,000 simulated subjects choosing at random on each line, the average
score was 0.60 and only 12 scored above 0.90. % ledger: b1:choi2007-ccei
Most people, given many budgets, choose as if they had stable preferences;
a minority do not.
\end{record}

\section*{Exercises}\label{sec:b1:preferences-and-utility:exercises}

\begin{exercise}[$\star$]\label{exo:b1:preferences-and-utility:mrs-cd}
Compute the \omterm{def:b1:preferences-and-utility:mrs}{MRS} of $u = x_1^{0.3}x_2^{0.7}$ at any bundle, then at $(1, 1)$
and at $(3, 7)$. What does each number mean?
\end{exercise}

\begin{exercise}[$\star$]\label{exo:b1:preferences-and-utility:axioms}
Which axiom (completeness, transitivity, monotonicity, convexity) does each
statement contradict?
\begin{enumerate}
  \item ``I prefer a flat near the station to one with a garden, one with a garden to one with a view, and one with a view to one near the station.''
  \item ``I cannot say whether I prefer a week in the mountains or a week at the sea, and I never will.''
  \item ``Give me all bread or all cheese, but not a bit of each.''
  \item ``Three loaves are worse than two, though I could give the third away for free.''
\end{enumerate}
\end{exercise}

\begin{exercise}[$\star$]\label{exo:b1:preferences-and-utility:read-chart}
From the opening chart's table, the indices for 1996 and 2024 are: all
items 71.74 and 126.07; rents 72.02 and 115.46; audio-visual equipment
380.88 and 71.64. % figdata: hicp_ratios
\begin{enumerate}
  \item By what factor did each change, and by what percentage?
  \item Did rents rise or fall relative to all items? By how much?
\end{enumerate}
\end{exercise}

\begin{exercise}[$\star$]\label{exo:b1:preferences-and-utility:log}
A consumer's preferences are represented by $u = x_1x_2$. Does
$v = \ln x_1 + \ln x_2$ represent the same preferences? And $w = -x_1x_2$?
And $z = (x_1x_2)^2$ on positive bundles? Justify each.
\end{exercise}

\begin{exercise}[$\star\star$]\label{exo:b1:preferences-and-utility:complements}
For $u = \min(x_1, x_2)$:
\begin{enumerate}
  \item draw three \omterm{def:b1:preferences-and-utility:indifference-curve}{indifference curves};
  \item show that at a bundle with $x_1 > x_2$ the \omterm{def:b1:preferences-and-utility:mrs}{MRS} is 0, and that at a bundle with $x_1 < x_2$ the \omterm{def:b1:preferences-and-utility:indifference-curve}{indifference curve} is vertical, so the \omterm{def:b1:preferences-and-utility:mrs}{MRS} is not a finite number;
  \item explain why no \omterm{def:b1:preferences-and-utility:mrs}{MRS} exists at a corner such as $(3, 3)$.
\end{enumerate}
\end{exercise}

\begin{exercise}[$\star\star$]\label{exo:b1:preferences-and-utility:quasi}
For $u = 4\sqrt{x_1} + x_2$, compute the \omterm{def:b1:preferences-and-utility:mrs}{MRS} at $(4, 1)$ and at $(4, 10)$.
Interpret the result when $x_2$ is money.
\end{exercise}

\begin{exercise}[$\star\star$]\label{exo:b1:preferences-and-utility:predict}
A consumer with $u = \sqrt{x_1x_2}$ buys electricity ($x_1$) and audio-visual
equipment ($x_2$), and chooses where the \omterm{def:b1:preferences-and-utility:mrs}{MRS} equals the price ratio
$p_1/p_2$ (\cref{prop:b1:constrained-optimisation:tangency}). Between 1996
and 2024 the price of electricity relative to equipment rose by a factor of
13.5. % figdata: hicp_ratios
\begin{enumerate}
  \item Predict the change in the ratio $x_2/x_1$ she chooses.
  \item Predict the change in the share of her spending on electricity.
  \item What observation in household budgets would show these predictions wrong?
\end{enumerate}
\end{exercise}

\begin{exercise}[$\star\star$]\label{exo:b1:preferences-and-utility:cycle}
A consumer is observed three times. At prices $p^1$ she buys $x^1$ although
$x^2$ was affordable; at $p^2$ she buys $x^2$ although $x^3$ was affordable;
at $p^3$ she buys $x^3$ although $x^1$ was affordable, and $x^1$ cost strictly
less than $x^3$ at $p^3$. Which axiom of the model of \omterm{def:b1:preferences-and-utility:rational}{rational preferences}, with
\omterm{def:b1:preferences-and-utility:monotone}{monotone} preferences and choice of the best affordable bundle, do these
choices contradict? Explain.
\end{exercise}

\begin{exercise}[$\star\star$]\label{exo:b1:preferences-and-utility:refute-convex}
A consumer, on every straight budget line she faces, spends all her money
on one good. Explain why this does not refute \omterm{def:b1:preferences-and-utility:convex}{convex preferences}. Then
describe two choices between pairs of bundles, using two bundles $x$, $z$,
their average $m = (x + z)/2$ and a slightly larger bundle $m^{+} > m$,
that would refute \omterm{def:b1:preferences-and-utility:rational}{rational}, \omterm{def:b1:preferences-and-utility:monotone}{monotone} and \omterm{def:b1:preferences-and-utility:convex}{convex preferences} together.
\end{exercise}

\begin{exercise}[$\star\star\star$]\label{exo:b1:preferences-and-utility:transform}
Let $u$ be differentiable with $u_2 > 0$ and $f$ differentiable with $f' > 0$.
Prove that $v = f(u)$ has the same \omterm{def:b1:preferences-and-utility:mrs}{MRS} as $u$ at every bundle, and give an
example of a transformation $f$ with $f' < 0$ somewhere for which $v$ does not
represent the same preferences.
\end{exercise}

\begin{exercise}[$\star\star\star$]\label{exo:b1:preferences-and-utility:convex-not-strict}
Show that the preferences represented by $u = \min(x_1, 2x_2)$ are \omterm{def:b1:preferences-and-utility:convex}{convex}
but not strictly \omterm{def:b1:preferences-and-utility:convex}{convex}: find two bundles $x \neq z$, both indifferent to
$y = (4, 2)$, whose average is also indifferent to $y$.
\end{exercise}

\begin{exercise}[$\star\star\star$]\label{exo:b1:preferences-and-utility:design}
Design two budget lines for two goods and a choice from each that
would reveal a cycle, so that no \omterm{def:b1:preferences-and-utility:rational}{rational}, \omterm{def:b1:preferences-and-utility:monotone}{monotone} preference ordering
could have produced them. Then say what the laboratory record in the
chapter found about how often real subjects make such choices.
\end{exercise}

\section{Problem: Two Flatmates and a Fridge}\label{sec:b1:preferences-and-utility:problem}

\begin{problem}[{Weekend problem --- the swap rate of milk for bread}]\label{pb:b1:preferences-and-utility:fridge}
Ana and Ben share a fridge in Brenmoor (invented). A bundle is written
(litres of milk, loaves of bread), $(m, b)$. Ana's preferences are represented by
$u_A = \sqrt{mb}$ and Ben's by $u_B = m + 6\ln b$. A \omterm{def:b1:preferences-and-utility:mrs}{marginal rate of
substitution} in this problem is in litres of milk per loaf:
$\mathrm{MRS} = (\partial u/\partial b)/(\partial u/\partial m)$. Four
bundles are on offer: $P = (4, 2)$, $Q = (2, 4)$, $R = (3, 3)$ and $S = (6, 1)$.

\medskip\noindent\textbf{Part I --- Rankings.}
\begin{enumerate}
  \item Rank the four bundles for Ana, using $\succ$ and $\sim$.
  \item Rank them for Ben (give his utility values to three decimals).
  \item A third flatmate, Cleo, says $P \succ Q$, $Q \succ R$ and $R \succ P$. Which axiom fails? Show how a fourth flatmate could gain from her by a sequence of swaps.
  \item $R$ is the average of $P$ and $Q$. Is Ana's ranking of $P$, $Q$, $R$ consistent with strictly \omterm{def:b1:preferences-and-utility:convex}{convex preferences}? And Ben's?
  \item Show that Ana's and Ben's preferences are \omterm{def:b1:preferences-and-utility:monotone}{monotone}.
  \item Can Ana and Ben disagree on a ranking and both be \omterm{def:b1:preferences-and-utility:rational}{rational}? Use your answers to questions 1 and 2.
\end{enumerate}

\medskip\noindent\textbf{Part II --- Utility and the MRS.}
\begin{enumerate}[resume]
  \item Show that $v_A = mb$ and $w_A = \ln m + \ln b$ represent Ana's preferences.
  \item Compute Ana's \omterm{def:b1:preferences-and-utility:mrs}{MRS} at any bundle, then at $P$ and at $Q$.
  \item Compute Ben's \omterm{def:b1:preferences-and-utility:mrs}{MRS} at any bundle, then at $P$ and at $Q$.
  \item Write Ana's \omterm{def:b1:preferences-and-utility:indifference-curve}{indifference curve} through $R$ as a function $m(b)$, and give its slope at $R$.
  \item Interpret Ben's \omterm{def:b1:preferences-and-utility:mrs}{MRS} at $Q$ in one sentence.
  \item Show that Ana's \omterm{def:b1:preferences-and-utility:mrs}{MRS} diminishes as $b$ rises along an \omterm{def:b1:preferences-and-utility:indifference-curve}{indifference curve}.
\end{enumerate}

\medskip\noindent\textbf{Part III --- Which form?}
\begin{enumerate}[resume]
  \item Ana's \omterm{def:b1:preferences-and-utility:mrs}{MRS} depends only on $m/b$. Which standard form is her utility, and with which exponent?
  \item Ben's \omterm{def:b1:preferences-and-utility:mrs}{MRS} does not depend on $m$. Which standard form is his?
  \item Dan uses milk only with bread, two litres per loaf, and wastes the rest. Write a utility for him and say what his \omterm{def:b1:preferences-and-utility:mrs}{MRS} is.
  \item Eva would always give three litres of milk for one loaf, or a loaf for three litres. Write a utility for her and give her \omterm{def:b1:preferences-and-utility:mrs}{MRS}.
  \item Eva holds $(9, 1)$ and may swap milk for bread at 1.4 litres per loaf as much as she likes. What does she choose?
  \item If every bundle of Ana's is scaled up by the same factor, what happens to her \omterm{def:b1:preferences-and-utility:mrs}{MRS}? And to Ben's?
\end{enumerate}

\medskip\noindent\textbf{Part IV --- A swap.}
Ana holds $(9, 1)$ and Ben holds $(1, 7)$.
\begin{enumerate}[resume]
  \item Compute each one's \omterm{def:b1:preferences-and-utility:mrs}{MRS} at these bundles. Who values bread more, at the margin?
  \item They swap one loaf for 2 litres of milk. Compute both utilities before and after. Do both gain?
  \item For a small swap of bread for milk at a rate of $r$ litres per loaf, for which rates do both gain?
  \item At a fixed rate $r$, Ana buys bread until her \omterm{def:b1:preferences-and-utility:mrs}{MRS} equals $r$, and Ben sells until his equals $r$. Show that Ana buys $(9 - r)/(2r)$ loaves and Ben sells $7 - 6/r$.
  \item Find the rate at which these two amounts are equal, and the swap.
  \item Give each one's bundle and utility after the swap, and both \omterm{def:b1:preferences-and-utility:mrs}{MRS}. State the result: the swap rate they agree on.
\end{enumerate}
\end{problem}
